\appendix
\section{Reproducción del proyecto y de las capturas}
\label{anx:repro}
\begin{enumerate}
	\item Instalar Flutter y Android Studio; comprobar con \texttt{flutter doctor}.
	\item Abrir una terminal en la carpeta \texttt{matrlab-flutter} y, solo la primera vez, ejecutar \texttt{flutter create --project-name matrlab --platforms android,web .} (agrega las carpetas de plataforma sin sobrescribir \texttt{lib/} ni \texttt{pubspec.yaml}).
	\item Ejecutar \texttt{.\textbackslash ejecutar.ps1}, que corre \texttt{flutter pub get}, \texttt{flutter analyze} y \texttt{flutter test}; guarda las 30 capturas en \texttt{capturas/} y las salidas en \texttt{salida\_analyze.txt}, \texttt{salida\_test.txt} e \texttt{historial\_git.txt}.
	\item \texttt{flutter run} abre la aplicación en un emulador; \texttt{flutter run -d chrome}, en el navegador.
	\item Credenciales de demostración: estudiante con un código de 8 dígitos (por ejemplo, \texttt{20201234}), cualquier clave de 4 caracteres o más y un ciclo (el 10 tiene el turno abierto; el 5 aún no); administrador con un código de 4 o más caracteres (por ejemplo, \texttt{admin}).
	\item Para ver la contingencia: ingresar como administrador, pestaña Operación, activar «Simular caída del servicio»; luego, como estudiante, intentar matricularse.
	\item Compilar este informe: \texttt{latexmk -pdf informe.tex} dentro de \texttt{informe\_lab05/}.
\end{enumerate}

\section{Registro de cambios del repositorio}
\label{anx:git}
El registro de la tabla~\ref{tab:git} se obtuvo con \texttt{git log}; las fechas son las reales de cada \emph{commit}. Los siete primeros pertenecen al Laboratorio~4; los demás, al Laboratorio~5.
\IfFileExists{historial.tex}{\begin{table}[H]
	\centering\footnotesize
	\begin{tabularx}{\textwidth}{|p{0.6cm}|p{1.5cm}|p{2.6cm}|p{0.9cm}|X|}
		\hline
		\cab\hd{N.°} & \hd{Commit} & \hd{Fecha} & \hd{Lab.} & \hd{Mensaje}\\
		\hline
		1 & \texttt{0012d28} & 25/10/2026 12:22 & 4 & Lab 04 Plataformas Emergentes: MatrLab en Flutter con informe LaTeX \\\hline
		2 & \texttt{c9281e4} & 25/10/2026 12:49 & 4 & Corrige iconos en las capturas de prueba y desborde del panel de administrador \\\hline
		3 & \texttt{f3bffc9} & 25/10/2026 12:52 & 4 & Corrige desborde en tarjetas de curso y estabiliza las pruebas de interfaz \\\hline
		4 & \texttt{dde5cc7} & 25/10/2026 12:56 & 4 & Ajusta la prueba de flujo: sin desplazamiento innecesario \\\hline
		5 & \texttt{4d47b27} & 25/10/2026 13:06 & 4 & Prueba de flujo: mismo ritmo de avance que las capturas \\\hline
		6 & \texttt{0faa6a1} & 25/10/2026 13:08 & 4 & Agrega capturas reales, resultados de las pruebas y registro de cambios 1.1 \\\hline
		7 & \texttt{da2c27c} & 25/10/2026 13:11 & 4 & Agrega la seccion de recorrido de la aplicacion solo con capturas \\\hline
		8 & \texttt{4110006} & 25/10/2026 13:31 & 4 & Initial commit \\\hline
		9 & \texttt{f010bf1} & 25/10/2026 13:36 & 4 & Merge branch 'main' of https://github.com/ahilacondo/matrlab-flutter \\\hline
		10 & \texttt{2b1cd0b} & 01/11/2026 14:07 & 5 & Lab 05: MatrLab 2.0 con turnos, atencion virtual y modo contingencia \\\hline
		11 & \texttt{41599fe} & 01/11/2026 14:07 & 5 & Agrega diagramas y mockups del Lab 05 \\\hline
		12 & \texttt{4a011ee} & 01/11/2026 14:08 & 5 & Agrega el informe LaTeX del Lab 05 (requisitos reales, Scrum y desarrollo) \\\hline
		13 & \texttt{97e950a} & 01/11/2026 15:23 & 5 & Rediseña figuras del informe del Lab 05 y pasa las tablas de Scrum a hojas de calculo \\\hline
		14 & \texttt{26ccd68} & 01/11/2026 15:23 & 5 & Merge remote-tracking branch 'origin/main' \\\hline
		15 & \texttt{3d29401} & 01/11/2026 15:31 & 5 & Corrige la prueba de contingencia: espera a que la hoja inferior termine de subir \\\hline
		16 & \texttt{023749d} & 01/11/2026 15:36 & 5 & Evita el simbolo de flecha en los turnos y usa el reloj fijo en los avisos iniciales \\\hline
	\end{tabularx}
	\caption{Registro de commits (\texttt{git log}), del más antiguo al más reciente, hasta el commit \texttt{023749d}, que es el estado probado. Los commits posteriores (informe y capturas finales) se registran en \texttt{historial\_git.txt}.}
	\label{tab:git}
\end{table}
}{\pendiente{Generar \texttt{historial.tex} con el registro de commits.}}
